\documentclass[11pt]{amsart}
\usepackage{amsmath,amssymb,amsthm}
\usepackage{booktabs}
\usepackage[margin=1.2in]{geometry}
\usepackage[colorlinks=true,linkcolor=blue,citecolor=blue,urlcolor=blue]{hyperref}
\usepackage{microtype}
\emergencystretch=1.5em
\raggedbottom

\newtheorem{theorem}{Theorem}[section]
\newtheorem{proposition}[theorem]{Proposition}
\newtheorem{corollary}[theorem]{Corollary}
\newtheorem{lemma}[theorem]{Lemma}
\newtheorem{conjecture}[theorem]{Conjecture}
\newtheorem{question}[theorem]{Question}
\theoremstyle{definition}
\newtheorem{definition}[theorem]{Definition}
\newtheorem{example}[theorem]{Example}
\theoremstyle{remark}
\newtheorem{remark}[theorem]{Remark}
\newtheorem{observation}[theorem]{Observation}

\newcommand{\CC}{\mathbb{C}}
\newcommand{\ZZ}{\mathbb{Z}}
\newcommand{\QQ}{\mathbb{Q}}
\newcommand{\NN}{\mathbb{N}}
\newcommand{\Crit}{\operatorname{Crit}}
\newcommand{\Bal}{\operatorname{Balanced}}
\newcommand{\rank}{\operatorname{rank}}
\newcommand{\Mbar}{\overline{\mathcal{M}}}
\newcommand{\gfrak}{\mathfrak{g}}
\newcommand{\chat}{\hat{c}}
\DeclareMathOperator{\Aut}{Aut}

\hyphenation{de-scen-dent de-scen-dents wall-cross-ing sin-gu-lar-i-ty}

\title[The central-charge threshold for wall-crossing in open FJRW theory]{The
central-charge threshold for wall-crossing\\ in two-variable open FJRW theory}

\author{Bernd Johannes Wuebben}
\subjclass[2020]{14N35 (primary); 14J33, 53D45, 14B05 (secondary)}
\keywords{Open FJRW theory, Landau--Ginzburg models, wall-crossing,
central charge, simple singularities, boundary conditions}
\date{July 19, 2026}

\begin{document}

\begin{abstract}
Gross, Kelly and Tessler constructed a genus-zero open enumerative theory for
the Landau--Ginzburg models $(x^{r}+y^{s},\mu_{r}\times\mu_{s})$, in which the
open FJRW invariants depend on a choice of boundary conditions and jump across
walls organized by a wall-crossing group. We determine exactly when this
dependence occurs. Primary wall-crossing is governed by the central charge: it
is absent precisely when $\chat_{W}<1$, that is, for the simple singularities
$A_{n-1}$, $D_{4}$, $E_{6}$, $E_{8}$, so the open theory detects the
classical ADE threshold. The minimal number of interior insertions
producing a wall obeys a uniform trichotomy; its smallest value, two
insertions, occurs first for $x^{4}+y^{4}$.
A single descendent insertion produces wall-crossing for every model. Each
wall-crossing generator is the Hamiltonian vector field, with respect to
$dx\wedge dy$, of the boundary monomial of its critical disk, the monomial
recording its boundary marked points; this exhibits the wall-crossing Lie
algebra as a graded Hamiltonian subalgebra with tropical-vertex structure
constants. Critical graphs of vanishing Euler weight (among them the unique
first-wall graph of each parabolic pair) generate the torus direction excluded
from the pro-nilpotent algebra, and it is there that the modularity of the
simply-elliptic singularities first registers in the open theory. The threshold statement agrees with, and for the
Fermat pairs gives an elementary proof of, a result established in greater
generality in Maher's thesis, of which we learned after this note was drafted.
The classification of first walls and the Hamiltonian normal form are the
starting point of the scattering-diagram construction of a companion paper.
All computations are exact and are reproduced by code accompanying the note.
\end{abstract}

\maketitle
\enlargethispage{3pt}

\section{Introduction}
\label{sec:intro}

\subsection*{Background}
For a nondegenerate quasi-homogeneous polynomial $W$ and a group $G$ of diagonal
symmetries, FJRW theory \cite{FJR13} assigns to the Landau--Ginzburg model
$(W,G)$ a cohomological field theory: a genus-zero $A$-model built from moduli
of $W$-spin curves and the Euler class of a Witten bundle. In their monograph
\cite{GKT}, Gross, Kelly and Tessler construct the \emph{open} analogue for the
two-variable Fermat model
\[
W=x^{r}+y^{s},\qquad G=\mu_{r}\times\mu_{s},\qquad r,s\geq 2,
\]
replacing closed $W$-spin curves by $W$-spin disks. The resulting open FJRW
invariants are integrals of multisections of a descendent Witten bundle over
moduli spaces that are real orbifolds with corners, and their central output is
an open mirror theorem: a generating function of the open invariants is a versal
deformation of $W$ whose oscillatory period integrals recover the closed
extended FJRW invariants, with the deformation parameters already flat
\cite[Theorem~0.3]{GKT}.

A structural feature distinguishes the open theory from its closed counterpart
and from the open $r$-spin theory of Buryak--Clader--Tessler
\cite{BCT21,BCT18}: the invariants are \emph{not} canonical. Boundary conditions
for the Witten bundle over the corners of moduli space are constrained, but not
determined, by the positivity requirements of \cite{GKT}; the residual freedom
makes the open invariants jump as the boundary data cross certain walls. Gross,
Kelly and Tessler organize these jumps into a pro-nilpotent Landau--Ginzburg
wall-crossing group $G^{r,s}$ acting on the set of possible systems of
invariants, and prove that this action is faithful and transitive
\cite[Theorems~0.7 and~4.26]{GKT}: the systems of open invariants form a single
torsor, while the period integrals -- hence the closed theory -- are constant
along it.

The source of every wall is a \emph{critical graph}: a stable $W$-spin
graph~$\Lambda$ whose Witten bundle has rank one above the dimension of its
moduli space, $\rank E_{\Lambda}(\mathbf d)=\dim\Mbar^{W}_{\Lambda}+1$
\cite[Definition~2.32]{GKT}. Such a graph indexes a one-parameter family of
codimension-one boundary strata through which zeros of the multisection can
escape. The wall-crossing Lie algebra $\gfrak^{r,s}$ is a direct sum of
one-dimensional pieces, one for each critical graph
\cite[Definition~4.22]{GKT}. Consequently the wall-crossing group is trivial --
and the invariants are canonical -- precisely when no critical graph exists.

\subsection*{The prediction}
Whether critical graphs exist depends on $(r,s)$ and on whether one allows
descendents. In the descendent-free (\emph{primary}) theory,
\cite[Remark~0.9]{GKT} predicts that wall-crossing is absent exactly for the
simple singularities, and asserts this for an explicit list of pairs. As stated,
that remark contains a discrepancy between its list and its naming, which we
record and correct below. The prediction itself -- that the primary open FJRW
invariants are canonical if and only if $W$ is a simple singularity -- is left
open in \cite{GKT} and is reiterated as such in the survey \cite{GKTsurvey}.

\subsection*{Results}
We prove the prediction and recast it as a clean invariant-theoretic threshold.
Recall that the \emph{central charge} of $W=x^{r}+y^{s}$, whose variables carry
weights $q_{x}=1/r$ and $q_{y}=1/s$, is
\begin{equation}
\label{eq:centralcharge}
\chat_{W}\;=\;(1-2q_{x})+(1-2q_{y})\;=\;2-\frac{2}{r}-\frac{2}{s}.
\end{equation}

\begin{theorem}[Central-charge criterion]
\label{thm:main}
Let $W=x^{r}+y^{s}$ with $r,s\geq2$ and central charge $\chat_{W}$ as
in~\eqref{eq:centralcharge}. The following are equivalent.
\begin{enumerate}
\item[(i)] $\chat_{W}<1$.
\item[(ii)] $\dfrac1r+\dfrac1s>\dfrac12$.
\item[(iii)] $W$ is a simple singularity: up to interchange of $x$ and $y$,
\[
(r,s)\in\{(n,2):n\geq2\}\cup\{(3,3),(3,4),(3,5)\},
\]
of type $A_{n-1}$, $D_{4}$, $E_{6}$, $E_{8}$ respectively.
\item[(iv)] For every finite marking set $I\subseteq\NN$, the genus-zero primary
open FJRW invariants of $(x^{r}+y^{s},\mu_{r}\times\mu_{s})$ are independent of
the choice of symmetric canonical boundary conditions; equivalently, the
descendent-free part of the wall-crossing group $G^{r,s}_{A_{I,\mathrm{sym}}}$
of \cite{GKT} (marginally enlarged as in Definition~\ref{rem:torusdir}) acts
trivially.
\end{enumerate}
When these fail, that is when $\chat_{W}\geq1$, there is a finite marking set
$I$ for which the primary invariants genuinely depend on the boundary
conditions.
\end{theorem}

\begin{definition}[Euler weight; the marginally enlarged group]
\label{rem:torusdir}
For interior data $(J,\mathbf d)$ set $w(J,\mathbf d)=m(J,\mathbf d)-rs$, the
\emph{Euler weight}, where $m(J,\mathbf d)$ is the mass of
\cite[Notation~2.29]{GKT} (its formula is recalled in
Proposition~\ref{prop:grading}); by the critical condition of
\cite[Proposition~2.37]{GKT}, a critical graph with data $(J,\mathbf d)$ and
index $(k_{1},k_{2})$ (the indexing of \eqref{eq:generator} below: the graph
has $k_{1}{+}1$ and $k_{2}{+}1$ boundary marked points of the two types)
satisfies $sk_{1}+rk_{2}=w(J,\mathbf d)$, so vanishing Euler weight means
index $(0,0)$. Such a \emph{marginal} critical graph has generator the torus
field $x\partial_{x}-y\partial_{y}$, the one direction that
\cite[Definition~4.22]{GKT} excludes for nilpotency of the filtration. The
\emph{marginally enlarged} wall-crossing Lie algebra and group are
\[
\widetilde{\gfrak}^{\,r,s}=\gfrak^{r,s}\ \oplus
\bigoplus_{w(J,\mathbf d)=0,\ r(J)=s(J)=0}
\QQ\,u_{J,\mathbf d}\,(x\partial_{x}-y\partial_{y}),
\qquad
\widetilde G^{r,s}=\exp\bigl(\widetilde{\gfrak}^{\,r,s}\bigr),
\]
with $u_{J,\mathbf d}=\prod_{i\in J}t_{a_{i},b_{i},d_{i}}$ the deformation
monomial; the exponential converges adically because every summand carries its
monomial. The enlargement is forced on us: the proofs of
\cite[Theorem~5.8 and Corollary~5.13]{GKT} use exactly these summands (their
jump sets impose only the critical condition, which admits the index
$(0,0)$), and their action realizes the explicit jumping of
\cite[Example~4.9]{GKTsurvey}; for the
elliptic cases this is the group appearing in \cite[\S5.2.2]{Maher}.
Faithfulness and transitivity extend: the enlarged group preserves the set of
systems of open invariants, acts faithfully and transitively on it, and
contains every jump automorphism of \cite[Theorem~5.8]{GKT}
(\cite[Lemma~2.4]{companion}, whose proof runs through the realizability step
(5.51)--(5.53) of \cite{GKT}). All statements below are over the enlarged
group.
\end{definition}

The equivalence (i)$\Leftrightarrow$(ii)$\Leftrightarrow$(iii) is elementary and
classical; the content is (iv). The threshold $\chat_{W}=1$ is exactly the
boundary between the simple singularities ($\chat_{W}<1$) and the parabolic
(simple-elliptic, $\chat_{W}=1$) and hyperbolic ($\chat_{W}>1$) regimes. Thus
the presence of primary wall-crossing detects, through open enumerative
geometry, the classical ADE threshold.

Beyond the dichotomy, the same count determines \emph{when} wall-crossing first
appears. For an interior insertion the relevant twists lie in the range
$0\leq a\leq r-2$, $0\leq b\leq s-2$; write $m$ for the number of interior
insertions (the mass $m(J,\mathbf d)$, always written with its arguments, is
a different quantity).

\begin{theorem}[First wall]
\label{thm:firstwall}
Let $m_{*}(r,s)$ be the least number of primary interior insertions admitting a
critical graph. Then
\[
m_{*}(r,s)=
\begin{cases}
\infty, & \tfrac1r+\tfrac1s>\tfrac12\quad(\text{simple}),\\[2pt]
2, & \min(r,s)\geq 4,\\[2pt]
3, & \min(r,s)=3\ \text{and}\ \max(r,s)\geq 6.
\end{cases}
\]
The case $m_{*}=2$ is first realized by $W=x^{4}+y^{4}$ with two interior
insertions of twist $(2,2)$, the example of \cite[\S6]{GKTsurvey}.
\end{theorem}

Finally, descendents change the picture entirely: they generate walls in every
case, including the simple singularities, as \cite[Remark~0.9]{GKT} anticipates.
The first such wall is universal.

\begin{proposition}[First descendent wall]
\label{prop:descendent}
For every $(r,s)$ with $r,s\geq2$, a single interior insertion carrying one
descendent is a critical graph. Hence the open FJRW invariants with descendents
exhibit wall-crossing for every two-variable Fermat model, simple or not.
\end{proposition}

Where walls do occur, their generators have a uniform Hamiltonian description.
Gross, Kelly and Tessler write the wall-crossing Lie algebra in terms of the
vector fields
$X_{a,b}=x^{a}y^{b}\bigl((b{+}1)x\partial_{x}-(a{+}1)y\partial_{y}\bigr)$
\cite[eqn.~(4.18)]{GKT} and observe that they preserve $dx\wedge dy$
\cite[p.~85]{GKT}. We name the Hamiltonian and read off the consequence.

\begin{proposition}[Hamiltonian wall generators]
\label{prop:ham}
For $\omega=dx\wedge dy$ and $H_{a,b}=x^{a+1}y^{b+1}$, the wall generator
$X_{a,b}$ is the Hamiltonian vector field of $H_{a,b}$; in particular it is
divergence-free. The assignment $H_{a,b}\mapsto X_{a,b}$ is an isomorphism of
graded Lie algebras from the Poisson algebra on the monomials $x^{p}y^{q}$,
$p,q\geq1$, onto its image, under which the bracket is the tropical-vertex
determinant
\[
[X_{a,b},X_{c,d}]=\bigl((b{+}1)(c{+}1)-(a{+}1)(d{+}1)\bigr)\,X_{a+c,\,b+d}.
\]
It restricts to an isomorphism onto the wall-crossing Lie algebra
$\mathfrak g^{r,s}$ of \cite{GKT}, which is thereby the graded Hamiltonian
subalgebra spanned by the critical Hamiltonians
$\{H_{\Lambda}=x^{k_{1}+1}y^{k_{2}+1}\}$, where the critical graph $\Lambda$ has
$k_{1}{+}1$ and $k_{2}{+}1$ boundary marked points of the two types, so that
$H_{\Lambda}$ is precisely the \emph{boundary monomial} of the critical disk.
The $(s,r)$-weighted degree of $H_{\Lambda}$ equals $w(\Lambda)+r+s$, with
$w$ the Euler weight of Definition~\ref{rem:torusdir}; marginal graphs
($H_{\Lambda}=xy$) contribute the adjoined torus direction of
Definition~\ref{rem:torusdir}. This makes the tropical-vertex analogy of
\cite[Remarks~0.8 and~4.19]{GKT} an explicit map.
\end{proposition}

Proposition~\ref{prop:ham} completes the program the criterion was part of: an
ADE/central-charge classification of when walls exist, a determination of the
first wall, and a Hamiltonian normal form for the wall generators. The one
structure it does \emph{not} supply --- a scattering diagram whose rays are the
critical graphs --- is constructed in the companion paper \cite{companion},
for which the census organized here is the starting point; see
\S\ref{sec:hamiltonian} and the closing section.

\subsection*{Method}
Nothing here rebuilds the machinery of \cite{GKT}. Theorem~\ref{thm:main}
follows by combining three of its results: the closed-form count of critical
graphs \cite[Notation~2.38]{GKT}, the identification of the wall-crossing Lie
algebra with the span of the critical graphs \cite[Definition~4.22]{GKT}, and
the faithfulness of the wall-crossing action \cite[Theorem~4.26]{GKT}. The
arithmetic core is a single inequality, Lemma~\ref{lem:bound}, bounding the
graph count by $m\,\chat_{W}$. All numerical claims (the ADE list, the trichotomy
$m_{*}(r,s)$, the first-wall census with its marginal/genuine split, and the
universal descendent wall) are proved in
\S\S\ref{sec:proof}--\ref{sec:parabolic} and independently confirmed by the
accompanying program \texttt{src/census.py}, which transcribes Notation~2.38
verbatim and checks these statements over a grid of $(r,s)$ using only integer
and rational arithmetic (\S\ref{sec:cert}).

\subsection*{Relation to prior work}
The uniqueness of the primitive form for simple singularities is classical
\cite{Saito83,NoumiYamada98}, and is the closed-mirror shadow of the statement
that these are the potentials without wall-crossing. On the open side,
Basalaev--Buryak constructed open Saito potentials for the $A$- and $D$-series
by solving the open WDVV equations \cite{BB21}; the present criterion explains,
uniformly, why the primary open theory is unambiguous on the whole ADE locus,
including the $E$-cases, for which no closed-form potential has yet been
computed. The wall-crossing group of \cite{GKT} is drawn, ``in spirit,'' by analogy
with the tropical vertex group of Kontsevich--Soibelman \cite{KS06} and its role
in Fano mirror symmetry \cite{Gross10} (\cite[Remark~0.8]{GKT}); separately, a
tropical version of open FJRW theory is raised as an open problem in
\cite[\S6]{GKTsurvey}, there through tropical moduli. The census of critical
graphs organized here is a prerequisite for a scattering-diagram form combining
the two, carried out in \cite{companion}. After this note was drafted we
learned that Maher's thesis establishes the threshold in substantially greater
generality, for invertible polynomials, via the (non)triviality of the
primary wall-crossing group \cite[Proposition~1.0.5]{Maher}; we keep the
present treatment for its elementary Fermat-case proof, the explicit first-wall
census, and the Hamiltonian normal form, on which the scattering construction
builds directly.

\section{The count of critical graphs}
\label{sec:count}

We recall only what is needed and refer to \cite{GKT} for all definitions. Fix
$W=x^{r}+y^{s}$. A genus-zero smooth $(r,s)$-disk carries interior marked points
with twists $(a_{i},b_{i})$, $0\leq a_{i}\leq r-1$, $0\leq b_{i}\leq s-1$, and a
descendent vector $\mathbf d=(d_{i})\in\ZZ_{\geq0}$. The interior insertions
contributing to the primary potential lie in the Neveu--Schwarz range
$0\leq a_{i}\leq r-2$, $0\leq b_{i}\leq s-2$ \cite[eqn.~(0.5)]{GKT}.

Given a set $J$ of interior markings with twists $(a_{i},b_{i})_{i\in J}$ and
descendents $\mathbf d=(d_{i})_{i\in J}$, define, following
\cite[Notation~2.29]{GKT},
\begin{equation}
\label{eq:ells}
\sum_{i\in J}a_{i}=r(J)+\ell_{1}r,\qquad
\sum_{i\in J}b_{i}=s(J)+\ell_{2}s,
\end{equation}
where $r(J)\in\{0,\dots,r-1\}$ and $s(J)\in\{0,\dots,s-1\}$ are the residues, so
that
\begin{equation}
\label{eq:elldef}
\ell_{1}=\Big\lfloor\tfrac{1}{r}\textstyle\sum_{i\in J}a_{i}\Big\rfloor,\qquad
\ell_{2}=\Big\lfloor\tfrac{1}{s}\textstyle\sum_{i\in J}b_{i}\Big\rfloor.
\end{equation}
The classification of balanced and critical graphs with interior markings
exactly $J$ is \cite[Notation~2.38]{GKT}: setting
\begin{equation}
\label{eq:N}
N(J,\mathbf d)\;=\;\ell_{1}+\ell_{2}-|J|+1+\sum_{i\in J}d_{i},
\end{equation}
one has $\bigl|\Bal(J,\mathbf d)\bigr|=N(J,\mathbf d)+1$ and
\begin{equation}
\label{eq:critcount}
\bigl|\Crit(J,\mathbf d)\bigr|=\max\bigl(N(J,\mathbf d),\,0\bigr).
\end{equation}
In particular a critical graph with interior markings $J$ and descendent vector
$\mathbf d$ exists if and only if $N(J,\mathbf d)\geq1$. (Throughout,
$J\neq\emptyset$: the boundary-only case $J=\emptyset$, where $N=1$, comprises
the two balanced, non-critical graphs of \cite[Example~3.4]{GKT}, and
\eqref{eq:critcount} is not asserted there.)

The reduction from invariants to graph counts is the following consequence of
\cite{GKT}, which we isolate for use in both directions of
Theorem~\ref{thm:main}.

\begin{proposition}
\label{prop:reduction}
Fix a finite marking set $I\subseteq\NN$ and work over the coefficient ring
$A_{I,\mathrm{sym}}=\QQ[t_{a,b,d}\,]/\mathrm{Ideal}(I)$ of \cite{GKT} (the
ideal caps each twist $(a,b)$ at the number of labels of $I$ carrying it). Let
$\mathbf d=\mathbf 0$.
\begin{enumerate}
\item[(a)] If no critical graph with descendent vector $\mathbf0$ has interior
markings contained in $I$, then $G^{r,s}_{A_{I,\mathrm{sym}}}$ fixes every
primary open FJRW invariant; the primary invariants are canonical.
\item[(b)] If some critical graph with descendent vector $\mathbf0$ has interior
markings contained in $I$, then two systems of symmetric canonical open FJRW
invariants over $I$ differ in a primary invariant.
\end{enumerate}
\end{proposition}

\begin{proof}
By \cite[Definition~4.22]{GKT} the wall-crossing Lie algebra
$\gfrak^{r,s}_{A_{I,\mathrm{sym}}}$ is the direct sum, over all critical graphs
$\Lambda$ with markings in $I$, of the one-dimensional spaces
$\QQ\cdot v_{\Lambda}$, where for a critical graph with interior markings $J$,
descendents $\mathbf d$, and boundary degrees $(k_{1},k_{2})$,
\begin{equation}
\label{eq:generator}
v_{\Lambda}=u_{J,\mathbf d}\,x^{k_{1}}y^{k_{2}}
\bigl((k_{2}+1)x\partial_{x}-(k_{1}+1)y\partial_{y}\bigr),\qquad
u_{J,\mathbf d}=\prod_{i\in J}t_{a_{i},b_{i},d_{i}} .
\end{equation}
Here $(k_{1},k_{2})$ is the index of \cite[Definition~4.22]{GKT}: the critical
graph has $k_{1}{+}1$ boundary marked points of type $\sigma_{1}$ and
$k_{2}{+}1$ of type $\sigma_{2}$ \cite[eqn.~(2.31), Proposition~2.37]{GKT}, and
the index satisfies the critical condition
$sk_{1}+rk_{2}=w(J,\mathbf d)$ of Definition~\ref{rem:torusdir}, matching the
constraints defining $\gfrak^{r,s}$; the case $(k_{1},k_{2})=(0,0)$ is the
marginal torus direction of Definition~\ref{rem:torusdir}. The group is
$G^{r,s}_{A_{I,\mathrm{sym}}}=\exp\bigl(\gfrak^{r,s}_{A_{I,\mathrm{sym}}}\bigr)$
and acts on the set $\Upsilon^{\mathrm{sym}}_{r,s}(I)$ of chamber indices
(\cite[Definition~3.46]{GKT}) by
$g(W^{\nu})=W^{g(\nu)}$ \cite[Theorem~4.26]{GKT}.

For (a): if there is no $\mathbf0$-critical graph, every generator
\eqref{eq:generator} has some $d_{i}>0$, hence $u_{J,\mathbf d}$ is divisible by
a variable $t_{a,b,d}$ with $d>0$. Acting on a potential
$W^{\nu}=x^{r}+y^{s}+\sum_{0\le a\le r-2,\ 0\le b\le s-2}t_{a,b,0}\,x^{a}y^{b}+\cdots$
\cite[eqn.~(0.5)]{GKT}, such a generator raises the total degree in the
positive-descendent variables and therefore leaves unchanged the coefficients of
the purely primary monomials $u_{J',\mathbf 0}\,x^{\alpha}y^{\beta}$. Since
each primary open invariant $\nu_{\Gamma,\mathbf 0}$, the coefficient attached
to a graph $\Gamma$ in \cite[eqn.~(0.4)]{GKT}, is exactly such a coefficient,
it is fixed by every generator, hence by $G^{r,s}_{A_{I,\mathrm{sym}}}$.

For (b): a $\mathbf 0$-critical graph $\Lambda$ furnishes a nonzero generator
$v_{\Lambda}$ with $u_{J,\mathbf 0}=\prod_{i\in J}t_{a_{i},b_{i},0}$ a monomial
in primary variables only. Fix any system
$\nu\in\Upsilon^{\mathrm{sym}}_{r,s}(I)$, a chamber index in the sense of
\cite[Definition~3.46]{GKT}; the set is nonempty by
\cite[Theorem~0.1]{GKT}. By faithfulness \cite[Theorem~4.26, Step~2]{GKT},
$\exp(v_{\Lambda})$ does not fix $\nu$; and to first order,
$\exp(v_{\Lambda})(W^{\nu})-W^{\nu}=v_{\Lambda}(x^{r}+y^{s})+\cdots$ has the term
\[
u_{J,\mathbf 0}\,\bigl(r(k_{2}+1)\,x^{k_{1}+r}y^{k_{2}}
-s(k_{1}+1)\,x^{k_{1}}y^{k_{2}+s}\bigr),
\]
a primary monomial with a nonzero coefficient. Thus $\nu$ and
$\exp(v_{\Lambda})\cdot\nu$ are two systems of symmetric canonical invariants
(both in the single orbit $\Upsilon^{\mathrm{sym}}_{r,s}(I)$, by transitivity)
that differ in a primary invariant.
\end{proof}

\section{Proof of the central-charge criterion}
\label{sec:proof}

The equivalences (i)$\Leftrightarrow$(ii)$\Leftrightarrow$(iii) are immediate:
\eqref{eq:centralcharge} gives $\chat_{W}<1\Leftrightarrow\frac2r+\frac2s>1$, and
$\frac1r+\frac1s>\frac12$ characterizes the two-variable Fermat simple
singularities, whose only types are $A_{n-1}=(n,2)$, $D_{4}=(3,3)$,
$E_{6}=(3,4)$ and $E_{8}=(3,5)$ (there is no Fermat $E_{7}$, as $E_{7}=x^{3}+xy^{3}$
is not of the form $x^{r}+y^{s}$; likewise $D_{n}=x^{n-1}+xy^{2}$ is Fermat only
for $n=4$, where $x^{3}+y^{3}\cong D_{4}$). It remains to prove
(i)$\Leftrightarrow$(iv), for which, by Proposition~\ref{prop:reduction}, it
suffices to decide the existence of a primary ($\mathbf d=\mathbf0$) critical
graph. The whole matter rests on one inequality.

\begin{lemma}
\label{lem:bound}
Let $J$ be a set of $m=|J|$ interior markings with Neveu--Schwarz twists,
$0\leq a_{i}\leq r-2$ and $0\leq b_{i}\leq s-2$. Then, with $\ell_{1},\ell_{2}$
as in \eqref{eq:elldef},
\[
\ell_{1}+\ell_{2}\;\leq\;m\,\chat_{W}.
\]
Consequently the primary count \eqref{eq:N} obeys
$N(J,\mathbf 0)\leq 1-m(1-\chat_{W})$.
\end{lemma}

\begin{proof}
Dropping the floor functions and using the twist bounds,
\[
\ell_{1}+\ell_{2}
=\Big\lfloor\tfrac1r\textstyle\sum a_{i}\Big\rfloor
+\Big\lfloor\tfrac1s\textstyle\sum b_{i}\Big\rfloor
\leq\frac{\sum a_{i}}{r}+\frac{\sum b_{i}}{s}
\leq\frac{m(r-2)}{r}+\frac{m(s-2)}{s}
=m\Big(2-\tfrac2r-\tfrac2s\Big)=m\,\chat_{W}.
\]
The second claim is $N(J,\mathbf 0)=\ell_{1}+\ell_{2}-m+1\leq m\chat_{W}-m+1$.
\end{proof}

\begin{proof}[Proof of (i)$\Rightarrow$(iv)]
Assume $\chat_{W}<1$. For any $J$ with $m=|J|\geq1$ interior markings,
Lemma~\ref{lem:bound} gives $\ell_{1}+\ell_{2}\leq m\chat_{W}<m$; as
$\ell_{1}+\ell_{2}$ is an integer strictly below the integer $m$, we get
$\ell_{1}+\ell_{2}\leq m-1$, i.e.\ $N(J,\mathbf 0)\leq0$. By
\eqref{eq:critcount} there is no primary critical graph, for any $J$ and any
finite $I$. Proposition~\ref{prop:reduction}(a) then says the primary invariants
are canonical for every $I$. (The case $m=0$ is vacuous: there is no interior
marking to twist, and $N=1$ corresponds to the two boundary-only balanced graphs
of \cite[Example~3.4]{GKT}, which are not critical.)
\end{proof}

\begin{proof}[Proof of $\neg$(i)$\Rightarrow\neg$(iv)]
Assume $\chat_{W}\geq1$, i.e.\ $\frac2r+\frac2s\leq1$; in particular
$r,s\geq3$ and not both equal to $3$, so $r,s\geq3$ with $\max(r,s)\geq4$, and
we may take all interior twists maximal, $(a_{i},b_{i})=(r-2,s-2)$. With $m$ such
markings,
\begin{equation}
\label{eq:maxN}
N=\Big\lfloor\tfrac{m(r-2)}{r}\Big\rfloor+\Big\lfloor\tfrac{m(s-2)}{s}\Big\rfloor-m+1
=\Big(m-\big\lceil\tfrac{2m}{r}\big\rceil\Big)
+\Big(m-\big\lceil\tfrac{2m}{s}\big\rceil\Big)-m+1
=m-\big\lceil\tfrac{2m}{r}\big\rceil-\big\lceil\tfrac{2m}{s}\big\rceil+1,
\end{equation}
using $\lfloor m(r-2)/r\rfloor=\lfloor m-2m/r\rfloor=m-\lceil2m/r\rceil$. We
exhibit $m$ with $N\geq1$, i.e.\ $\lceil2m/r\rceil+\lceil2m/s\rceil\leq m$.

If $\frac2r+\frac2s<1$ (the hyperbolic regime $\chat_{W}>1$), then
$\lceil2m/r\rceil+\lceil2m/s\rceil\leq\frac{2m}{r}+\frac{2m}{s}+2
=m\bigl(\frac2r+\frac2s\bigr)+2$, which is $\leq m$ once
$m\bigl(1-\frac2r-\frac2s\bigr)\geq2$; such $m$ exists since
$1-\frac2r-\frac2s>0$.

If $\frac2r+\frac2s=1$ (the parabolic regime $\chat_{W}=1$: the pairs $(4,4)$,
$(3,6)$, $(6,3)$), choose $m$ with $r\mid2m$ and $s\mid2m$, e.g.\
$m=\operatorname{lcm}(r,s)/2$, an integer for all three parabolic pairs; then
$\lceil2m/r\rceil+\lceil2m/s\rceil=\frac{2m}{r}+\frac{2m}{s}=m$,
so $N=1$.

In either case a primary critical graph exists once $I$ contains $m$ interior
labels of maximal twist, which is possible because each twist is realized
infinitely often in the marking data \cite[\S2.4]{GKT}. By
Proposition~\ref{prop:reduction}(b) the primary invariants over such an $I$
depend on the boundary conditions. This proves $\neg$(iv), and with it
Theorem~\ref{thm:main}.
\end{proof}

\begin{remark}[On Remark~0.9 of \cite{GKT}]
\label{rem:erratum}
\cite[Remark~0.9]{GKT} states that the pairs
$(r,s)\in\{(n,2),(2,n),(3,3),(3,4),(4,3)\}$ give simple singularities
``$A_{n-1}$, $E_{6}$, and $E_{8}$''. Two adjustments reconcile the list with the
naming. First, $(3,3)$ appears in the list but $D_{4}$ is not named; $x^{3}+y^{3}$
is the $D_{4}$ singularity (three distinct lines, Milnor number $4$). Second,
$E_{8}$ is named but its pairs $(3,5),(5,3)$ are absent from the list; these are
$x^{3}+y^{5}$, Milnor number $8$. The corrected wall-crossing-free set is
\[
\{(n,2),(2,n):n\geq2\}\cup\{(3,3),(3,4),(4,3),(3,5),(5,3)\},
\]
of types $A_{n-1},D_{4},E_{6},E_{8}$, exactly as in Theorem~\ref{thm:main}(iii).
The theorem shows this list is not merely a set of examples but the complete
answer to the prediction.
\end{remark}

\section{The first wall and the descendent wall}
\label{sec:firstwall}

\begin{proof}[Proof of Theorem~\ref{thm:firstwall}]
The simple case $m_{*}=\infty$ is Theorem~\ref{thm:main}. Assume $\chat_{W}\geq1$
and use \eqref{eq:maxN}, which gives the maximal $N$ over size-$m$ primary
multisets, since $N$ is nondecreasing in the twists and maximal at
$(r-2,s-2)$. Thus $m_{*}(r,s)$ is the least $m$ with
$\lceil2m/r\rceil+\lceil2m/s\rceil\leq m$.

\emph{$\min(r,s)\geq4$.} For $m=1$, $\lceil2/r\rceil+\lceil2/s\rceil=1+1=2>1$
(each $2/r,2/s\leq\frac12$). For $m=2$ and $r,s\geq4$,
$\lceil4/r\rceil+\lceil4/s\rceil=1+1=2=m$, so $N=1$ and $m_{*}=2$.

\emph{$\min(r,s)=3$, say $r=3$, with $\max=s\geq6$.} Here $a_{i}\leq1$, so
$\ell_{1}=\lfloor m/3\rfloor$. For $m=2$,
$\lceil4/3\rceil+\lceil4/s\rceil=2+1=3>2$. For $m=3$,
$\lceil6/3\rceil+\lceil6/s\rceil=2+1=3=m$ (using $s\geq6$), so $N=1$ and
$m_{*}=3$. For $3\le s\le5$ the pair is simple and $m_{*}=\infty$, already
covered.

Finally $\min(r,s)=2$ forces the $A$-type simple case, $m_{*}=\infty$. This
exhausts all $(r,s)$. For $x^{4}+y^{4}$, the witnessing multiset is two
insertions of twist $(r-2,s-2)=(2,2)$, giving $N=1$: the example of
\cite[\S6]{GKTsurvey}.
\end{proof}

\begin{proof}[Proof of Proposition~\ref{prop:descendent}]
Take $|J|=1$ with any twist. Then $\ell_{1}=\ell_{2}=0$ and, by \eqref{eq:N},
$N(J,\mathbf d)=0+0-1+1+d=d$, where $d$ is the descendent degree of the single
insertion. Hence $d=1$ gives $N=1$: exactly one critical graph, for every
$(r,s)$. Proposition~\ref{prop:reduction}(b) is stated for
$\mathbf d=\mathbf 0$, but its proof uses only that the generator is nonzero
and that the group acts faithfully: applied to the descendent-$1$ generator,
the same first-order computation produces a monomial divisible by $t_{a,b,1}$
with nonzero coefficient, so two systems of invariants differ in a descendent
invariant. Hence the descendent theory wall-crosses in every case.
\end{proof}

\begin{remark}
The trichotomy of Theorem~\ref{thm:firstwall} shows the $x^{4}+y^{4}$ example of
\cite{GKTsurvey} is not isolated: two interior insertions already force a wall
for every $(r,s)$ with $r,s\geq4$, and three suffice whenever
$\min(r,s)=3<\frac12\max(r,s)$. Combined with
Proposition~\ref{prop:descendent}, the qualitative picture is complete: the
descendent-free theory is canonical exactly on the ADE locus, while switching on
a single descendent breaks canonicity everywhere.
\end{remark}

\section{The wall generators are Hamiltonian}
\label{sec:hamiltonian}

We now prove Proposition~\ref{prop:ham}. Throughout, $\omega=dx\wedge dy$, and
for a bivariate polynomial $H$ the Hamiltonian vector field is
$X_{H}:=(\partial_{y}H)\,\partial_{x}-(\partial_{x}H)\,\partial_{y}$, so that
$\iota_{X_{H}}\omega=dH$. In the rank-two case the wall-crossing Lie algebra
$\mathfrak g^{r,s}$ of \cite{GKT} is a sum of one-dimensional spaces spanned by
\begin{equation}
\label{eq:Xab}
X_{a,b}=x^{a}y^{b}\bigl((b+1)x\partial_{x}-(a+1)y\partial_{y}\bigr)
=(b+1)x^{a+1}y^{b}\,\partial_{x}-(a+1)x^{a}y^{b+1}\,\partial_{y},
\end{equation}
one for each critical graph, the exponents $(a,b)=(k_{1},k_{2})$ being its
\cite[Definition~4.22]{GKT} index: the graph has $k_{1}{+}1$ and $k_{2}{+}1$
boundary marked points of the two types, so the Hamiltonian
$H_{k_{1},k_{2}}=x^{k_{1}+1}y^{k_{2}+1}$ of Lemma~\ref{lem:hamvf} below is its
boundary monomial.

\begin{lemma}
\label{lem:hamvf}
For all $a,b\geq0$, $X_{a,b}=X_{H_{a,b}}$ with $H_{a,b}=x^{a+1}y^{b+1}$. In
particular $\operatorname{div}X_{a,b}=0$, so $X_{a,b}$ preserves $\omega$.
\end{lemma}

\begin{proof}
$\partial_{y}H_{a,b}=(b+1)x^{a+1}y^{b}$ and
$\partial_{x}H_{a,b}=(a+1)x^{a}y^{b+1}$, so
$X_{H_{a,b}}=(b+1)x^{a+1}y^{b}\partial_{x}-(a+1)x^{a}y^{b+1}\partial_{y}$, which is
\eqref{eq:Xab}. A Hamiltonian vector field is divergence-free:
$\operatorname{div}X_{a,b}=\partial_{x}\partial_{y}H_{a,b}-\partial_{y}\partial_{x}H_{a,b}=0$,
recovering the volume-preservation of \cite[p.~85]{GKT}.
\end{proof}

\begin{lemma}
\label{lem:bracket}
For all $a,b,c,d\geq0$,
\begin{equation}
\label{eq:bracketid}
[X_{a,b},X_{c,d}]=\bigl((b+1)(c+1)-(a+1)(d+1)\bigr)\,X_{a+c,\,b+d}.
\end{equation}
Equivalently, with the Poisson bracket $\{H,K\}:=X_{H}(K)$ induced by $\omega$,
\[
\{H_{a,b},H_{c,d}\}=\bigl((b+1)(c+1)-(a+1)(d+1)\bigr)\,H_{a+c,\,b+d}.
\]
\end{lemma}

\begin{proof}
Write $X_{a,b}=f\partial_{x}+g\partial_{y}$ with $f=(b+1)x^{a+1}y^{b}$,
$g=-(a+1)x^{a}y^{b+1}$, and $X_{c,d}=F\partial_{x}+G\partial_{y}$ similarly. A
direct computation of
$[X_{a,b},X_{c,d}]=(X_{a,b}F-X_{c,d}f)\partial_{x}+(X_{a,b}G-X_{c,d}g)\partial_{y}$
gives, for the $\partial_{x}$-coefficient,
\[
X_{a,b}F-X_{c,d}f
=(b+d+1)\bigl[(b+1)(c+1)-(a+1)(d+1)\bigr]\,x^{a+c+1}y^{b+d},
\]
which is $\bigl((b+1)(c+1)-(a+1)(d+1)\bigr)$ times the $\partial_{x}$-coefficient
$(b+d+1)x^{a+c+1}y^{b+d}$ of $X_{a+c,b+d}$; the $\partial_{y}$-coefficient matches
identically. For the Poisson form,
$\{H_{a,b},H_{c,d}\}=X_{a,b}(H_{c,d})
=(b+1)(c+1)x^{a+c+1}y^{b+d+1}-(a+1)(d+1)x^{a+c+1}y^{b+d+1}$, and
$X_{H_{a+c,b+d}}=X_{a+c,b+d}$ by Lemma~\ref{lem:hamvf} shows the two identities
agree.
\end{proof}

The determinant $\langle(a,b),(c,d)\rangle=(b+1)(c+1)-(a+1)(d+1)$ appearing
in~\eqref{eq:bracketid} is the structure constant of the tropical vertex Lie
algebra of Kontsevich--Soibelman \cite{KS06}; see \cite{Gross10}. Assembling
Lemmas~\ref{lem:hamvf} and~\ref{lem:bracket} proves Proposition~\ref{prop:ham}:
$H_{a,b}\mapsto X_{a,b}$ is a Lie-algebra isomorphism from the Poisson algebra on
$\{x^{p}y^{q}:p,q\geq1\}$ onto its image, restricting to an isomorphism onto
$\mathfrak g^{r,s}$ once one keeps only the Hamiltonians $H_{\Lambda}$ of critical
graphs. The grading is immediate from the critical condition.

\begin{proposition}
\label{prop:grading}
Let $\Lambda$ be a critical graph with interior markings $J$, descendents
$\mathbf d$, and index $(k_{1},k_{2})$, and set
$m(\Lambda)=m(J,\mathbf d)=rs+\sum_{i\in J}(sa_{i}+rb_{i}+rs(d_{i}-1))$
\cite[Notation~2.29]{GKT}. Then the Hamiltonian $H_{\Lambda}=x^{k_{1}+1}y^{k_{2}+1}$
has $(s,r)$-weighted degree
\[
s(k_{1}+1)+r(k_{2}+1)=m(\Lambda)+r+s-rs .
\]
Thus $\mathfrak g^{r,s}$ is graded by the $(s,r)$-weighted degree of the wall
Hamiltonian, and the degree records the mass of the critical graph
(equivalently its Euler weight $w(\Lambda)=m(\Lambda)-rs$ of
Definition~\ref{rem:torusdir}).
\end{proposition}

\begin{proof}
The critical condition is $sk_{1}+rk_{2}=m(\Lambda)-rs$
\cite[Proposition~2.37]{GKT}; add $r+s$.
\end{proof}

\begin{remark}[Relation to the tropical vertex algebra]
\label{rem:tropical}
The tropical vertex Lie algebra is the Hamiltonian algebra of the \emph{torus}
form $d\log x\wedge d\log y$, for which
$\{x^{u},x^{v}\}=\langle u,v\rangle\,x^{u+v}$ with $\langle u,v\rangle$ the
determinant, so that Hamiltonians add exponents. Proposition~\ref{prop:ham} uses
instead the \emph{affine-plane} form $dx\wedge dy=xy\,(d\log x\wedge d\log y)$,
under which $\{x^{u},x^{v}\}=\langle u,v\rangle\,x^{u+v-(1,1)}$. The two Poisson
structures, hence the two Lie algebras, are intertwined by the exponent shift
$u\mapsto u-(1,1)$ induced by this factor $xy$ -- the same change of primitive
form that makes $dx\wedge dy$ the natural volume in \cite[Theorem~0.3]{GKT}. So
$\mathfrak g^{r,s}$ is precisely an affine-plane avatar of a tropical-vertex Lie
algebra, cut out by the critical Hamiltonians, which makes explicit the analogy
that \cite[Remarks~0.8 and~4.19]{GKT} draw ``in spirit'' with the
Kontsevich--Soibelman group. The genuine scattering diagram whose walls are the
critical graphs is constructed in \cite{companion}: there the set of invariant
systems is presented as the chamber set of a consistent scattering structure, a
canonical diagram with computable wall functions is extracted, and the
potentials are shown to be broken-line transports of a canonical seed system.
The tropical open FJRW theory called for in \cite[\S6]{GKTsurvey} is posed
there through tropical moduli, a priori a distinct route to the same tropical
picture.
\end{remark}

\section{The first wall at the parabolic boundary}
\label{sec:parabolic}

The threshold $\chat_{W}=1$ is not only where wall-crossing begins across the
family; the walls it produces have a uniform shape. By the \emph{first wall}
of $(r,s)$ we mean the set of critical graphs over interior multisets of the
minimal size $m_{*}(r,s)$. Combining the count
\eqref{eq:critcount} with the Hamiltonian dictionary of \S\ref{sec:hamiltonian}
gives a complete census of the first wall for each $(r,s)$; at the
parabolic boundary it collapses to a single wall.

\begin{proposition}[Parabolic first wall]
\label{prop:parabolic}
For any interior multiset $J$ with $N(J,\mathbf 0)=1$, the index of its unique
critical graph is $(k_{1},k_{2})=(r(J),s(J))$, of Euler weight
$s\,r(J)+r\,s(J)$. For the
three parabolic pairs, $(r,s)\in\{(4,4),(3,6),(6,3)\}$ with $\chat_{W}=1$, the
first wall consists of exactly one critical graph, its multiset forced to be
the maximal-twist one: two copies of $(2,2)$ for $(4,4)$, three copies of
$(1,4)$ for $(3,6)$, three copies of $(4,1)$ for $(6,3)$. Its residues vanish,
so its index is $(0,0)$, its weight is zero, and its direction is the
\emph{marginal torus direction} $x\partial_{x}-y\partial_{y}$ of
Definition~\ref{rem:torusdir}, with Hamiltonian $xy$. By contrast, for the
hyperbolic pairs of Table~\ref{tab:census} the first wall comprises several
critical graphs, both genuine positive-weight walls and marginal torus
directions.
\end{proposition}

\begin{proof}
If $N(J,\mathbf 0)=1$ there is a unique critical graph and, by
\cite[Notation~2.38, eqn.~(2.31)]{GKT} read in the indexing of
\eqref{eq:generator}, its index is $(k_{1},k_{2})=(r(J),s(J))$; the weight
statement is the critical condition of Definition~\ref{rem:torusdir}. For each parabolic pair
the minimal critical multiset is uniquely the maximal-twist one: at $m=m_{*}$
the inequality $\ell_{1}+\ell_{2}\geq m$ is attained only with equality and
only at $(a_{i},b_{i})=(r-2,s-2)$ (for $(4,4)$, $\ell_{1},\ell_{2}\leq1$ force
$a_{i}=b_{i}=2$; for $(3,6)$, $\ell_{1}\leq1,\ell_{2}\leq2$ force
$a_{i}=1,b_{i}=4$). There the residues vanish, so $(k_{1},k_{2})=(0,0)$: the
torus direction, with Hamiltonian $H_{0,0}=xy$ (Lemma~\ref{lem:hamvf}). The
hyperbolic data is the direct enumeration recorded in Table~\ref{tab:census}.
\end{proof}

\begin{table}[ht]
\centering
\begin{tabular}{ccccccl}
\toprule
$(r,s)$ & $\chat_{W}$ & regime & $m_{*}$ & graphs & marginal & wall Hamiltonians\\
\midrule
$(4,4)$ & $1$ & parabolic & $2$ & $1$ & $1$ & --- (torus $xy$ only)\\
$(3,6)$ & $1$ & parabolic & $3$ & $1$ & $1$ & --- (torus $xy$ only)\\
$(6,3)$ & $1$ & parabolic & $3$ & $1$ & $1$ & --- (torus $xy$ only)\\
$(4,5)$ & $11/10$ & hyperbolic & $2$ & $2$ & $1$ & $xy^{2}$\\
$(3,7)$ & $22/21$ & hyperbolic & $3$ & $2$ & $1$ & $xy^{2}$\\
$(4,6)$ & $7/6$ & hyperbolic & $2$ & $4$ & $2$ & $xy^{2},\,xy^{3}$\\
$(5,5)$ & $6/5$ & hyperbolic & $2$ & $5$ & $2$ & $xy^{2},\,x^{2}y,\,x^{2}y^{2}$\\
\bottomrule
\end{tabular}
\caption{The first wall of $x^{r}+y^{s}$: the number of critical graphs at
$m=m_{*}(r,s)$, how many are marginal (index $(0,0)$, the torus direction with
Hamiltonian $xy$), and the boundary monomials $x^{k_{1}+1}y^{k_{2}+1}$
Hamiltonian to the genuine wall generators. At the parabolic boundary the
first wall is purely marginal; genuine walls appear only in the hyperbolic
regime. Computed exactly; see \S\ref{sec:cert}.}
\label{tab:census}
\end{table}

The parabolic pairs are the two-variable Fermat \emph{simply-elliptic}
singularities ($x^{4}+y^{4}=\widetilde E_{7}$, $x^{3}+y^{6}=\widetilde E_{8}$),
for which the primitive form ceases to be $dx\wedge dy$ and becomes a period of
an elliptic curve; the marginal torus direction is where that transition first
registers in the open theory: the elliptic wall-crossing group of
\cite[\S5.2.2]{Maher} is exactly this torus factor, and the explicit jumping of
the $x^{4}+y^{4}$ invariants in \cite[Example~4.9]{GKTsurvey} is its action.
Genuine (pro-nilpotent) walls appear precisely in the hyperbolic regime. The
first \emph{symmetric} hyperbolic Fermat pair is $x^{5}+y^{5}$, with boundary
monomials $xy^{2},x^{2}y,x^{2}y^{2}$ in Table~\ref{tab:census}; its wall
functions and scattering diagram are computed in \cite{companion}.

\section{Verification}
\label{sec:cert}

The combinatorial and algebraic claims are verified exactly by the
accompanying program \texttt{src/census.py}, published with this note and the
companion paper in the repository of \cite{companion}, which uses only integer,
\texttt{fractions.Fraction}, and exact symbolic arithmetic. It transcribes
\eqref{eq:N} from \cite[Notation~2.38]{GKT} and checks the following.
\begin{itemize}
\item Over the grid $2\leq r,s\leq40$: the threefold identity of
Theorem~\ref{thm:main}(i)--(iii), and that the pairs without a primary
critical graph are exactly the ADE Fermat list, with $D_{4}=(3,3)$ and
$E_{8}=(3,5),(5,3)$ present (Remark~\ref{rem:erratum}). For $\chat_{W}<1$
the negative is proved rather than merely searched, Lemma~\ref{lem:bound}
bounding the search by $m\geq1/(1-\chat_{W})$, beyond which $N\leq0$ is
forced.
\item The trichotomy for $m_{*}(r,s)$ of Theorem~\ref{thm:firstwall}.
\item The first-wall census of Table~\ref{tab:census} with its
marginal/genuine split, each positive-weight index cross-checked against the
independent implementation of \cite[Definition~4.22]{GKT} in
\texttt{src/gkt\_algebra.py} of \cite{companion}, and the parabolic statement
of Proposition~\ref{prop:parabolic}.
\item The universal descendent wall of Proposition~\ref{prop:descendent},
including its refinement (torus direction at the zero twist, genuine wall of
index $(a,b)$ otherwise).
\item By exact polynomial arithmetic on vector fields: the Hamiltonian
identity of Lemma~\ref{lem:hamvf}, the divergence-freeness, and the bracket
\eqref{eq:bracketid} of Lemma~\ref{lem:bracket} over $0\leq a,b,c,d\leq4$.
\end{itemize}

\section*{Concluding remarks and open directions}
\addcontentsline{toc}{section}{Concluding remarks and open directions}

Of three natural continuations, the first, promoting the Hamiltonian
identification to a genuine scattering diagram whose walls are the critical
graphs, is carried out in the companion paper \cite{companion}: the set of open FJRW invariant systems is presented there as
the chamber set of a consistent scattering structure over $\mathfrak g^{r,s}$,
a canonical diagram $\mathfrak D^{r,s}$ with computable wall functions is
constructed from unique extreme chamber indices, the open topological
recursion appears as the direction transverse to the walls, and the potential
family is shown to be the broken-line transport of a canonical seed system.
The $(1,1)$-shift of Remark~\ref{rem:tropical} is precisely the mechanism that
forbids backscattering of broken lines there. Two directions remain open.

\begin{enumerate}
\item \emph{Explicit ADE potentials.} On the wall-crossing-free locus the
primary potential is unambiguous; computing it in closed form for $D_{4}$,
$E_{6}$, $E_{8}$ (matching the $D$-series potentials of \cite{BB21} and
extending them to the $E$-cases) and relating the coefficients to
Coxeter/root-system data is a concrete computational program.
\item \emph{The parabolic boundary.} Proposition~\ref{prop:parabolic} pins the
first wall of the parabolic pairs to the marginal torus direction. What
remains is the analytic side: these are the simply-elliptic singularities
$\widetilde E_{7},\widetilde E_{8}$, whose primitive form is a period of an
elliptic curve, so the flat structure is governed by modular forms rather than
polynomials \cite[\S4.2.2]{Maher}. Identifying the modular period map on the
torus factor, beyond the $x^{4}+y^{4}$ jumping of
\cite[Example~4.9]{GKTsurvey}, is the first genuinely transcendental case of
the theory.
\end{enumerate}

\begin{thebibliography}{GKT26}

\bibitem[BB21]{BB21} A.~Basalaev, A.~Buryak,
\emph{Open Saito theory for $A$ and $D$ singularities},
Int.\ Math.\ Res.\ Not.\ 2021, no.~7, 5460--5491. arXiv:1909.00598.

\bibitem[BCT22]{BCT18} A.~Buryak, E.~Clader, R.~J.~Tessler,
\emph{Open $r$-spin theory I: Foundations},
Int.\ Math.\ Res.\ Not.\ 2022, no.~14, 10458--10532. arXiv:2003.01082.

\bibitem[BCT24]{BCT21} A.~Buryak, E.~Clader, R.~J.~Tessler,
\emph{Open $r$-spin theory II: The analogue of Witten's conjecture for $r$-spin
disks}, J.\ Differential Geom.\ 128 (2024), no.~1, 1--75. arXiv:1809.02536.

\bibitem[FJR13]{FJR13} H.~Fan, T.~Jarvis, Y.~Ruan,
\emph{The Witten equation, mirror symmetry, and quantum singularity theory},
Ann.\ of Math.\ 178 (2013), no.~1, 1--106.

\bibitem[GKT22]{GKT} M.~Gross, T.~L.~Kelly, R.~J.~Tessler,
\emph{Open FJRW theory and mirror symmetry},
arXiv:2203.02435v2 (2022).

\bibitem[GKT26]{GKTsurvey} M.~Gross, T.~L.~Kelly, R.~J.~Tessler,
\emph{Open enumerative geometries for Landau--Ginzburg models},
arXiv:2602.12707v1 (2026).

\bibitem[Gro10]{Gross10} M.~Gross,
\emph{Mirror symmetry for $\mathbb{P}^{2}$ and tropical geometry},
Adv.\ Math.\ 224 (2010), no.~1, 169--245.

\bibitem[KS06]{KS06} M.~Kontsevich, Y.~Soibelman,
\emph{Affine structures and non-Archimedean analytic spaces},
in \emph{The unity of mathematics}, Progr.\ Math.\ 244, Birkh\"auser (2006),
321--385.

\bibitem[Mah24]{Maher} R.~Maher,
\emph{Predictions in open Fan--Jarvis--Ruan--Witten theory via mirror symmetry,
modularity, and wall-crossing}, Ph.D.\ thesis, University of Birmingham, 2024.
\texttt{etheses.bham.ac.uk/id/eprint/15556}.

\bibitem[NY98]{NoumiYamada98} M.~Noumi, Y.~Yamada,
\emph{Notes on the flat structures associated with simple and simply elliptic
singularities}, in \emph{Integrable systems and algebraic geometry}, World Sci.\
(1998), 373--383.

\bibitem[Sai83]{Saito83} K.~Saito,
\emph{Period mapping associated to a primitive form},
Publ.\ Res.\ Inst.\ Math.\ Sci.\ 19 (1983), no.~3, 1231--1264.

\bibitem[Wue26]{companion} B.~J.~Wuebben,
\emph{The scattering structure of open FJRW theory of $x^{r}+y^{s}$},
paper and accompanying computations,
\url{https://github.com/bwuebben/open-fjrw-scattering} (2026).

\end{thebibliography}

\bigskip
\noindent{\sc Bernd J. Wuebben, New York, NY,} \texttt{wuebben@gmail.com}

\end{document}
